% 08-conclusion-future-work.tex
% Section 8 — Conclusion and future work. Drafted G3.7. [ERIC-REVIEW] on
% Phase Legendre description and openness discussion.

\section{Conclusion and Future Work}
\label{sec:conclusion}

\subsection{Summary}

We have presented a Lean~4 formalisation of a conditional non-existence
result for nontrivial odd Collatz cycles (Theorem~\ref{thm:main}), with
an explicit, auditable separation between \emph{published} hypotheses
(Baker~\cite{Baker1966}, Barina~\cite{Barina2025}) and a \emph{derived}
hypothesis (\textsf{DerivedLargeKBound}). The formalisation contains zero
\texttt{sorry}, and the axiom chain of the main theorem consists of the
three standard kernel axioms of Mathlib and no others. An end-to-end
reproducibility pipeline, a continuous integration workflow, and an
internal audit trail under \texttt{docs/BIBLE/} accompany the paper and
the repository.

\subsection{Phase Legendre : formalising \textsf{DerivedLargeKBound}}

% [ERIC-REVIEW] Phase Legendre description and timeline. Please confirm
% the plan below is accurate and that 2--3 months is a realistic
% estimate.

The main item of planned future work is to promote
\textsf{DerivedLargeKBound} from a structure to a proven Lean theorem,
thereby removing it from the list of conditional hypotheses. We refer to
this follow-up as \emph{Phase Legendre}, in reference to Legendre's
best-approximation theorem on convergents of continued fractions, which
is the central analytic tool involved.

The plan has three parts :

\begin{enumerate}
  \item \emph{Mathlib survey.} Determine the current state of the
        \texttt{Mathlib.NumberTheory.ContinuedFractions} subsystem and
        identify which pieces (convergents of \(\log_2 3\), best-
        approximation properties, effective estimates) are already
        available and which need to be added.
  \item \emph{Formal derivation.} Formalise the argument sketched in
        Section~\ref{sec:hypotheses}, which combines Hypothesis~\ref{hyp:baker}
        with a convergent-based windowing argument, to produce a Lean
        proof of the predicate currently held by
        \textsf{DerivedLargeKBound}.
  \item \emph{Integration.} Replace the \texttt{structure} at the
        definition site in \texttt{Phase59ContinuedFractions.lean} by the
        newly proven theorem, keeping the statement of
        Theorem~\ref{thm:main} unchanged.
\end{enumerate}

Upon completion of Phase Legendre, the axiom chain of
Theorem~\ref{thm:main} will additionally include
\texttt{Lean.ofReduceBool} and \texttt{Lean.trustCompiler}, introduced by
the \texttt{native\_decide} tactic used in the arithmetic gap lemmas
(\(\mathtt{cf\_gap\_*}\), \(\mathtt{cf\_nbound\_*}\)) that underlie the
derivation. This change will be declared preemptively in
\texttt{expected\_axioms.md} before the integration commit and
acknowledged in a subsequent version of this paper.

A rough estimate for Phase Legendre is two to three additional months of
work, subject to a survey of Mathlib's current continued-fractions
infrastructure.

\subsection{Openness and limitations}

The present formalisation leaves several avenues open :

\begin{itemize}
  \item \textbf{Tightening the threshold.} The specific value \(k = 1322\)
        in Hypothesis~\ref{hyp:derived} is a consequence of taking
        \(\mu = 6\) rather than the finer Rhin bound. A refined
        formulation with \(\mu = 5.125\) would lower the threshold and
        could shorten the case split; this is a candidate for a
        follow-up technical note.
  \item \textbf{Auditing more theorems.} The probe
        \texttt{probes/check\_central\_axioms.lean} currently covers
        ten theorems (seven central, three illustrative
        \texttt{native\_decide} auxiliaries). Extending the baseline to
        all \(\mathtt{cf\_gap\_*}\) and \(\mathtt{cf\_nbound\_*}\)
        variants would remove residual ambiguity for reviewers
        consulting the code in depth.
  \item \textbf{Interoperating with other formalisations.} As other
        Lean formalisations of partial Collatz results
        appear~(see references in Section~\ref{sec:introduction}),
        reusing common lemmas (e.g.\ Syracuse maps, 2-adic valuations,
        Steiner-style equations) across repositories would be mutually
        beneficial. Such interoperation is not part of the present
        scope.
\end{itemize}

\subsection{Closing remark}

% [ERIC-REVIEW] Closing phrasing — please amend to your voice. Currently
% reserved and conditional, matching the rest of the paper.

We view this work as an intermediate step towards a fully unconditional
formalisation of the no-nontrivial-odd-cycle result, not as a final
statement on the subject. Its value, we hope, lies in the explicit
articulation of the conditional form, in the reproducibility and audit
infrastructure that accompanies it, and in the concrete path it identifies
for the remaining Phase Legendre work.
